\section{Cycles}
\label{section-cycles}

\noindent
Since we are not assuming our schemes are quasi-compact we have
to be a little careful when defining cycles. We have to allow
infinite sums because a rational function may have infinitely many
poles for example. In any case, if $X$ is quasi-compact then a
cycle is a finite sum as usual.

\begin{definition}
\label{definition-cycles}
Let $(S, \delta)$ be as in Situation \ref{situation-setup}.
Let $X$ be locally of finite type over $S$.
Let $k \in \mathbf{Z}$.
\begin{enumerate}
\item A {\it cycle on $X$} is a formal sum
$$
\alpha = \sum n_Z [Z]
$$
where the sum is over integral closed subschemes $Z \subset X$,
each $n_Z \in \mathbf{Z}$, and the collection
$\{Z; n_Z \not = 0\}$ is locally finite
(Topology, Definition \ref{topology-definition-locally-finite}).
\item A {\it $k$-cycle} on $X$ is a cycle
$$
\alpha = \sum n_Z [Z]
$$
where $n_Z \not = 0 \Rightarrow \dim_\delta(Z) = k$.
\item The abelian group of all $k$-cycles on $X$ is denoted $Z_k(X)$.
\end{enumerate}
\end{definition}

\noindent
In other words, a $k$-cycle on $X$
is a locally finite formal $\mathbf{Z}$-linear
combination of integral closed subschemes of $\delta$-dimension $k$.
Addition of $k$-cycles $\alpha = \sum n_Z[Z]$ and
$\beta = \sum m_Z[Z]$ is given by
$$
\alpha + \beta = \sum (n_Z + m_Z)[Z],
$$
i.e., by adding the coefficients.

\begin{remark}
\label{remark-cycles-pointwise}
Let $(S, \delta)$ be as in Situation \ref{situation-setup}.
Let $X$ be locally of finite type over $S$. Let $k \in \mathbf{Z}$.
Then we can write
$$
Z_k(X) = \bigoplus\nolimits_{\delta(x) = k}' K_0^M(\kappa(x))
\quad\subset\quad
\bigoplus\nolimits_{\delta(x) = k} K_0^M(\kappa(x))
$$
with the following notation and conventions:
\begin{enumerate}
\item $K_0^M(\kappa(x)) = \mathbf{Z}$ is the degree $0$ part of
the Milnor K-theory of the residue field $\kappa(x)$ of the point
$x \in X$ (see Remark \ref{remark-gersten-complex-milnor}), and
\item the direct sum on the right is over all points $x \in X$
with $\delta(x) = k$,
\item the notation $\bigoplus'_x$ signifies that we consider the
subgroup consisting of locally finite elements; namely, elements
$\sum_x n_x$ such that for every quasi-compact open $U \subset X$
the set of $x \in U$ with $n_x \not = 0$ is finite.
\end{enumerate}
\end{remark}

\begin{definition}
\label{definition-support-cycle}
Let $(S, \delta)$ be as in Situation \ref{situation-setup}.
Let $X$ be locally of finite type over $S$.
The {\it support} of a cycle $\alpha = \sum n_Z [Z]$ on $X$
is
$$
\text{Supp}(\alpha) = \bigcup\nolimits_{n_Z \not = 0} Z \subset X
$$
\end{definition}

\noindent
Since the collection $\{Z; n_Z \not = 0\}$ is locally finite
we see that $\text{Supp}(\alpha)$ is a closed subset of $X$.
If $\alpha$ is a $k$-cycle, then every irreducible component $Z$
of $\text{Supp}(\alpha)$ has $\delta$-dimension $k$.

\begin{definition}
\label{definition-effective-cycle}
Let $(S, \delta)$ be as in Situation \ref{situation-setup}.
Let $X$ be locally of finite type over $S$.
A cycle $\alpha$ on $X$ is {\it effective} if it
can be written as $\alpha =\sum n_Z [Z]$ with $n_Z \geq 0$ for all $Z$.
\end{definition}

\noindent
The set of all effective cycles is a monoid because the sum of
two effective cycles is effective, but it is not a group
(unless $X = \emptyset$).




